% GENERATED FILE. Edit canonical lessons, then run npm run generate:books.
% canonical-source-hash: fnv1a64:ec01fc8c90b60489
% canonical-lessons: TE-C264-cekku, TE-C264-visa, TE-C264-samanu, TE-C264-vartapatrika, TE-C264-nambaru

\chapter{Cheque, Visa, Luggage, Newspaper, Number}
\label{ch:te-cheque-visa-luggage-newspaper-number}
\hlchaptermodality{264}

\begin{chapteropening}
\textbf{By the end of this chapter:} \emph{I can say a cheque, a visa, luggage, a newspaper and a number.}

Complete the last lesson of chapter 264: I can say a cheque, a visa, luggage, a newspaper and a number.
\end{chapteropening}

\section[\texorpdfstring{\te{చెక్కు} (cekku)}{చెక్కు (cekku)}]{\te{చెక్కు} (cekku) — a cheque}

\label{lesson:TE-C264-cekku}



\begin{quote}
Before the new one: say the Telugu for shorts, then the Telugu for a vest.
\end{quote}

\subsection*{You'll want to know: \te{చెక్కు}}
\textbf{\te{చెక్కు}} — \emph{cekku} — \textquotedblleft{}a cheque\textquotedblright{}.

\begin{grammarlens}[title={What you've built}]
One more word for this chapter.
\end{grammarlens}

\subsection*{Your turn}
\begin{itemize}
\raggedright
  \item \emph{Say it:} \emph{cekku}
  \item \emph{Say it:} \emph{cekku}, once more
  \item \emph{Say it:} say \emph{cekku}
  \item \emph{Recall:} say the Telugu for shorts, then the Telugu for a vest, then say \emph{cekku} again
\end{itemize}

\begin{tcolorbox}[breakable,colback=teal!4,colframe=teal!35!black,title={Before you move on}]
What does \te{చెక్కు} mean? (\textquotedblleft{}A cheque\textquotedblright{}.) Say it once more.
\end{tcolorbox}
\section[\texorpdfstring{\te{వీసా} (vīsā)}{వీసా (vīsā)}]{\te{వీసా} (vīsā) — a visa}

\label{lesson:TE-C264-visa}



\begin{quote}
Before the new one: say the Telugu for a vest, then the Telugu for a cheque.
\end{quote}

\subsection*{You'll want to know: \te{వీసా}}
\textbf{\te{వీసా}} — \emph{vīsā} — \textquotedblleft{}a visa\textquotedblright{}.

\begin{grammarlens}[title={What you've built}]
One more word for this chapter.
\end{grammarlens}

\subsection*{Your turn}
\begin{itemize}
\raggedright
  \item \emph{Say it:} \emph{vīsā}
  \item \emph{Say it:} \emph{vīsā}, once more
  \item \emph{Say it:} say \emph{vīsā}
  \item \emph{Recall:} say the Telugu for a vest, then the Telugu for a cheque, then say \emph{vīsā} again
\end{itemize}

\begin{tcolorbox}[breakable,colback=teal!4,colframe=teal!35!black,title={Before you move on}]
What does \te{వీసా} mean? (\textquotedblleft{}A visa\textquotedblright{}.) Say it once more.
\end{tcolorbox}
\section[\texorpdfstring{\te{సామాను} (sāmānu)}{సామాను (sāmānu)}]{\te{సామాను} (sāmānu) — luggage, belongings}

\label{lesson:TE-C264-samanu}



\begin{quote}
Before the new one: say the Telugu for a cheque, then the Telugu for a visa.
\end{quote}

\subsection*{You'll want to know: \te{సామాను}}
\textbf{\te{సామాను}} — \emph{sāmānu} — \textquotedblleft{}luggage, belongings\textquotedblright{}.

\begin{grammarlens}[title={What you've built}]
One more word for this chapter.
\end{grammarlens}

\subsection*{Your turn}
\begin{itemize}
\raggedright
  \item \emph{Say it:} \emph{sāmānu}
  \item \emph{Say it:} \emph{sāmānu}, once more
  \item \emph{Say it:} say \emph{sāmānu}
  \item \emph{Recall:} say the Telugu for a cheque, then the Telugu for a visa, then say \emph{sāmānu} again
\end{itemize}

\begin{tcolorbox}[breakable,colback=teal!4,colframe=teal!35!black,title={Before you move on}]
What does \te{సామాను} mean? (\textquotedblleft{}Luggage, belongings\textquotedblright{}.) Say it once more.
\end{tcolorbox}
\section[\texorpdfstring{\te{వార్తాపత్రిక} (vārtāpatrika)}{వార్తాపత్రిక (vārtāpatrika)}]{\te{వార్తాపత్రిక} (vārtāpatrika) — a newspaper}

\label{lesson:TE-C264-vartapatrika}



\begin{quote}
Before the new one: say the Telugu for a visa, then the Telugu for luggage.
\end{quote}

\subsection*{You'll want to know: \te{వార్తాపత్రిక}}
\textbf{\te{వార్తాపత్రిక}} — \emph{vārtāpatrika} — \textquotedblleft{}a newspaper\textquotedblright{}.

\begin{grammarlens}[title={What you've built}]
One more word for this chapter.
\end{grammarlens}

\subsection*{Your turn}
\begin{itemize}
\raggedright
  \item \emph{Say it:} \emph{vārtāpatrika}
  \item \emph{Say it:} \emph{vārtāpatrika}, once more
  \item \emph{Say it:} say \emph{vārtāpatrika}
  \item \emph{Recall:} say the Telugu for a visa, then the Telugu for luggage, then say \emph{vārtāpatrika} again
\end{itemize}

\begin{tcolorbox}[breakable,colback=teal!4,colframe=teal!35!black,title={Before you move on}]
What does \te{వార్తాపత్రిక} mean? (\textquotedblleft{}A newspaper\textquotedblright{}.) Say it once more.
\end{tcolorbox}
\section[\texorpdfstring{\te{నంబరు} (nambaru)}{నంబరు (nambaru)}]{\te{నంబరు} (nambaru) — a number (phone, house)}

\label{lesson:TE-C264-nambaru}



\begin{quote}
Before the new one: say the Telugu for luggage, then the Telugu for a newspaper.
\end{quote}

\subsection*{You'll want to know: \te{నంబరు}}
\textbf{\te{నంబరు}} — \emph{nambaru} — \textquotedblleft{}a number (phone, house)\textquotedblright{}.

\begin{grammarlens}[title={What you've built}]
One more word for this chapter.
\end{grammarlens}

\subsection*{Your turn}
\begin{itemize}
\raggedright
  \item \emph{Say it:} \emph{nambaru}
  \item \emph{Say it:} \emph{nambaru}, once more
  \item \emph{Say it:} say \emph{nambaru}
  \item \emph{Recall:} say the Telugu for luggage, then the Telugu for a newspaper, then say \emph{nambaru} again
\end{itemize}

\begin{tcolorbox}[breakable,colback=teal!4,colframe=teal!35!black,title={Before you move on}]
What does \te{నంబరు} mean? (\textquotedblleft{}A number (phone, house)\textquotedblright{}.) Say it once more.
\end{tcolorbox}
